% ===== §2 The structure of the phenomenon =====
\section{The Formal Structure of the Datum}
\label{sec:structure}

\noindent
The description of \xref{sec:phenomenology} is now analysed for its structure, within the descriptive attitude and without explanatory supplement. The objective is to exhibit the form already present in the datum, so that the ontology of \xref{sec:ontology} is constrained by a determinate structure rather than by an impression. The three features identified above compose a single structure, the components of which are mutually dependent in the manner set out below.

\subsection{Simultaneity as constitutive medium}
\label{sec:simultaneity}

The togetherness of the co-apprehension is not the synchrony of two independently occurring events. On the construal as synchrony, each party would possess a distinct apprehension, the two apprehensions coinciding in time by attunement or chance. That construal is inadequate to the datum. In the structure as described, there are not two apprehensions that coincide but one apprehension that is irreducibly of two parties. The togetherness is not a relation obtaining between two antecedent awarenesses; it is the medium within which each awareness is first constituted. The apprehension of one party is internally determined by that party's apprehension of the other's apprehension, and conversely, such that neither awareness is well-defined independently of the other. This is the structure the prior treatment located at the constitution of the subject, where the first and second persons are mutually constituted rather than antecedently given and subsequently joined \citep{buber1970}; here the same structure obtains at the constitution of a value between two parties. The structure recurs at a higher level: in \pinkref{Paper~V} joint attention constitutes the persons, and here it constitutes the value the constituted persons sustain between them. The recurrence of the founding structure at successive levels is the first instance of a pattern the paper thematises in \xref{sec:breaking}.

\subsection{Mutual witness as reflexive and constitutive}
\label{sec:witness}

The reciprocal awareness is reflexive: what each party apprehends includes the other's apprehension, and what is apprehended includes that reciprocal inclusion. The resulting regress, each party apprehends the term, apprehends the other apprehending it, apprehends the other apprehending its own apprehension, is not to be truncated but is the structural form of the density the description recorded. The reflexive witnessing is not a spectatorial relation to a completed fact; it is a component of the term's generation. Its function is constitutive: were the witnessing withdrawn, were one party to simulate apprehension while in fact absent, the term would not be generated, however exact the simulation. The structure is therefore not counterfeitable from a single side. This yields, at the level of description, the condition the ethics of \xref{sec:keystone} discharges: the value subsists only in a witnessing that is genuinely reciprocal, and is extinguished when one party becomes a spectator of, or an extractor from, the other's investment.

\subsection{Directedness upon an opening}
\label{sec:opening}

The co-apprehension is directed, it has the intentional form of a directedness-upon, but its term is an opening, a site at which something is not yet determinate, rather than an object. The psychoanalytic tradition supplies an account of an attention organised around a void without filling it; the contemplative traditions supply several. The prior paper established the corresponding thesis in its own terms: real value, the value capable of endogenously bearing the positive phase of a good cycle, is the value organised around an absence and resistant to symbolisation. The present analysis adds that this organisation is, at its origin, joint: the void is co-encircled, attended by two parties conjointly. The nameless term is therefore not an object the parties have not yet named but an opening the parties jointly sustain. A consequence follows directly. To name the term, to fill the opening with a settled sign, terminates the generation; this is the pathology the prior paper identified as settling failure. The subsistence of the structure consists in the non-termination of the term, the sustaining of the opening, which is the condition the practice of \xref{sec:praxis} undertakes to tend.

\subsection{The constitutive identity of apprehension and generation}
\label{sec:hinge}

The three features converge on the structural result that organises the paper. The co-apprehension of the nameless value is identically its generation; the epistemic operation and the ontological event are one. Conjoined with the third feature, this yields a structure of consequence. Because the directedness is upon an opening rather than an object, there is no antecedent value for the apprehension to copy; and because apprehension is generation, the value generated is conditioned by the manner of the apprehending. The manner of attending, the void around which, the register within which, the quality with which the opening is sustained, determines the value generated. There is no fact, antecedent to and independent of the joint apprehending, as to what the value is; the apprehending is not the measurement of an antecedent quantity but a participation that brings the quantity to determinacy. This is the structural antecedent of the formal model of \xref{sec:grammar}, where it receives the articulation of a participatory measurement whose basis is selected in the manner of attending.

\subsection{The structure stated}
\label{sec:gathered}

The structure is stated as a single result, to be discharged by the ontology of the following section.

\begin{claim}[The structure of co-apprehended value]
In the founding datum, value is co-generated in co-apprehension, with four components: (i)~the apprehension is irreducibly of two parties, the togetherness constituting the medium within which each awareness arises rather than a relation between antecedent awarenesses; (ii)~the reciprocal witnessing is reflexive and constitutive, a component of the term's generation rather than a report upon a completed fact; (iii)~the joint attention is directed upon a sustained opening rather than upon an object, such that settling it with a sign terminates it; and (iv)~apprehension and generation are one operation, whence the manner of the joint apprehending conditions the value generated. It follows that value, so structured, is not an entity antecedent to its co-apprehension, and that co-apprehension is not a representation of antecedent states: the datum determines a non-objectual ontology of value and a participatory epistemology of its apprehension, antecedently to the introduction of any formal apparatus.
\end{claim}

\noindent
This result is the constraint on what follows. The phenomenological analysis determines both the being of value and the mode of its cognition, in a determinate direction: away from the object and from representation, toward generation and participation. The following section discharges the first constraint, the ontological, by establishing what relation, grammar, and value must be in order that the datum of \xref{sec:phenomenology} be possible. The formal epistemologies are admitted only after that determination is made.
